\section{符号说明}
\begin{longtable}{p{0.2\linewidth}p{0.68\linewidth}}
\caption{主要符号及含义}\label{tab:q1-notation}\\
\toprule 符号 & 含义 \\
\midrule\endfirsthead
\toprule 符号 & 含义 \\
\midrule\endhead
$\Omega$ & 半径 $1800\,\mathrm m$ 的目标圆域 \\
$g,p_i$ & 源位置、 第 $i$ 个检测位置 \\
$\theta_i,\varepsilon$ & 示向度、测向误差界 \\
$W_i$ & 第 $i$ 次测量的误差楔形 \\
$\mathcal P_k$ & 频道 $k$ 当前可能位置集合 \\
$D(\mathcal P)$ & 可能区域直径 \\
$(c^*,R^*)$ & 可能区域的最小包围圆圆心和半径 \\
$L,N_{\rm sw},N_{\rm mea}$ & 总路程、换频次数、测量次数 \\
$N_{\rm clr},N_{\rm suc}$ & 清除尝试次数、成功次数 \\
$U(q,g)$ & 候选测点 $q$ 在假定位置 $g$ 下的线性化误差代理 \\
$\lambda$ & 移动距离权重 \\
$r_o,r_i$ & 问题四外环、内环半径 \\
$n_s,n_m,n_c,n_f$ & 站点、测量、清除和光学兜底动作计数 \\
$T_{\rm move},T_{\rm mea},T_{\rm clr}$ & 移动、测量和清除对应的虚拟时间分量 \\
\bottomrule
\end{longtable}

表~\ref{tab:q1-notation} 中的符号贯穿四个子问题；其中 $\mathcal P_k$ 和 $R^*$ 是从测量反馈到安全清除判定的主要状态量。
